\documentclass[10pt]{article}
\usepackage[margin=0.82in]{geometry}
\usepackage{amsmath,amssymb}
\usepackage[T1]{fontenc}
\usepackage[utf8]{inputenc}
\usepackage[colorlinks=true,linkcolor=blue,urlcolor=blue,citecolor=blue]{hyperref}

\title{The Blackadar--Handelman Choquet-Simplex Conjecture\\in Stable Rank One}
\author{\texttt{fa\_banach\_001}, \texttt{agent\_lane\_12}}
\date{13 August 2026}

\begin{document}
\maketitle

\section*{Status}

\textbf{Literature-implied answer, partial subcase.}
Conjecture 1.2 in De Silva's arXiv:1601.03475 is fully settled for every
unital $C^*$-algebra of stable rank one by Antoine--Perera--Robert--Thiel,
arXiv:1809.03984, Theorem 4.1.  The later authors identify the classical
Blackadar--Handelman conjecture, but do not identify De Silva's paper; the
cross-paper implication is therefore recorded as literature-implied.

\section*{Source conjecture}

On PDF page 1 of \emph{A note on two Conjectures on Dimension functions of
$C^*$-algebras}, Conjecture 1.2 states that the affine space $DF(A)$ of
normalized dimension functions is a Choquet simplex for every $C^*$-algebra
$A$.  The source's Theorem 3.3 verifies this only when the extreme boundary of
the quasitrace simplex is finite and either the radius of comparison is finite
or $W(A)$ is almost unperforated.

\section*{Later theorem}

Antoine, Perera, Robert, and Thiel prove in Theorem 4.1 of
\emph{$C^*$-algebras of stable rank one and their Cuntz semigroups}
(PDF page 13 of the current arXiv version):

\begin{quote}
If $A$ is a unital $C^*$-algebra of stable rank one, then $K_0^*(A)$ is an
interpolation group and $DF(A)$ is a Choquet simplex.
\end{quote}

The proof first establishes Riesz interpolation for $\mathrm{Cu}(A)$.  Stable
rank one makes the classical Cuntz semigroup $W(A)$ hereditary in
$\mathrm{Cu}(A)$, so interpolation passes to $W(A)$ and then to its
Grothendieck group $K_0^*(A)$.  The normalized state space of an interpolation
group is a Choquet simplex, and this state space is naturally $DF(A)$.

Thus Conjecture 1.2 holds throughout the stable-rank-one class with no
simplicity, separability, exactness, finite-radius, or finite-extreme-boundary
assumption.

\section*{Scope}

This theorem does not settle the source's Conjecture 1.1 concerning density of
lower semicontinuous dimension functions, and it does not settle Conjecture 1.2
for arbitrary stable rank.  Recent primary-source surveys continue to state
the unrestricted Choquet-simplex problem as a conjecture.  A separate attempt
note records eight routes explored toward a full proof or counterexample and
the interpolation/realizability obstructions encountered.

\begin{thebibliography}{9}
\raggedright
\bibitem{DeSilva}
K. De Silva,
\emph{A note on two Conjectures on Dimension functions of $C^*$-algebras},
arXiv:1601.03475.

\bibitem{ABPT}
R. Antoine, F. Perera, L. Robert, and H. Thiel,
\emph{$C^*$-algebras of stable rank one and their Cuntz semigroups},
Duke Math. J. 171 (2022), 33--99;
arXiv:1809.03984.
\end{thebibliography}

\end{document}
